% !TEX program = xelatex
% 上海电机学院 毕业设计（论文）开题报告 示例（文字取自学校模板）
% 编译：latexmk proposal.tex
\documentclass{sdjuproposal}
\addbibresource{refs-proposal.bib}

% ---------------- 封面信息 ----------------
\title{兆瓦级风力发电机控制系统设计}
\subtitle{液压系统}            % 无副标题可删去
\school{电气学院}
\major{电气工程及其自动化}
\classname{XXXXXX}
\studentid{XXXXXXXXX}
\author{XXX}
\supervisor{XXX}
\finaldate{20XX\hspace{0.5em}年\hspace{0.5em}XX\hspace{0.5em}月\hspace{0.5em}XX\hspace{0.5em}日}

\keywords{风力发电，液压系统，制动控制，变桨距控制，PLC}

\begin{document}

\makecover

\maketitle
\begin{abstract}
风能作为一种可再生能源，己经成为一种新兴的重要发电型式。本文对兆瓦级风力发电机组电控系统中的液压系统进行研究和设计。并针对兆瓦级水平轴风力发电机的控制要求利用PLC来完成液压系统的制动控制和变桨距控制。
\end{abstract}

\section{文献综述}

随着全球能源短缺、环境污染以及温室效应等\cite{belyaev2005}问题的日益突显。寻找可持续的能源道路，开发利用新能源无疑是出路之一。目前现有的可再生能源有太阳能、地热能、风能、海洋能、生物质能、氢能和水能\cite{wang2006}。而在这些可再生能源之中，风能的利用（即风力发电）无疑是世界上增长最快的可再生能源\cite{long2003}。

我国是世界主要的煤炭生产和消耗国，常规能源危机及其造成的污染后果已日渐引起人们的重视，而储量丰富的可再生能源正尚待开发。尤其是新疆、内蒙、甘肃一带的丰富风资源既可缓解当地电力资源的缺乏，又可促进生产力与经济的发展。据统计，仅新疆九大风区的风能资源年蕴藏量即达到8093亿千瓦时，可装风机8100万千瓦，年发电量2449亿千瓦时，是新疆目前全部发电量的20倍，可见风力发电前景十分广阔，开发利用风能必将成为二十一世纪的重要议程\cite{li2005,gong2005}。

风力发电的工作原理是利用风轮将风能转变为机械能，然后利用发电机将机械能转变为电能\cite{sun1997}。风力发电机由风轮、机舱和塔架组成\cite{wu2006}。机舱作为风力发电机控制系统的核心。其主要又由偏航系统、齿轮箱、联轴器、安全离合器和液压系统组成。而在机舱中最为重要的控制系统则是液压系统\cite{liu2004,yang2006}。……

\section{选题背景及其意义}

节约能源，提高能源利用率，大力开发使用新能源和可再生能源，逐步以洁净能源替代矿物燃料，是我国能源建设与发展应遵循的原则，也是实施可持续发展战略的一个重要组成部分，对于环境保护和增加能源供应有着积极作用。风力发电做为无污染的可再生能源随着世界范围内石油、煤炭储量的不断减少和燃用石油、煤炭对环境污染的严重影响，使它越来越受到人们的关注，风力机制造技术的不断进步和单机容量不断增大，使风电成本逐年下降，目前已达到接近火电成本并可与水电和核电相竞争的水平，具有十分显著的效益和发展前景。但我国的风力发电机大多引进国外整套设备，从中国大范围、持久开发风能的需要来看，单纯依赖国外进口风机绝不是根本出路。只有在引进国外先进技术的同时发展我们自己的风机制造业，才是百年大计。因此研制具有自主知识产权的风力发电机具有十分重大的意义。

本课题中风力发电机液压控制系统主要是，………………

\section{研究内容}

\subsection{液压系统的基本工作原理}
依靠密封容积的……………

\subsection{液压泵}
液压泵是液压系统的动力元件，它是将……………

\subsection{XXXX}
……。

\subsection{XXXX}
……

\subsection{XXXX}
……

\section{工作特色及其难点，拟采取的解决措施}

\subsection{工作特色及其难点}
随着风力发电技术的不断发展，风力发电机的单机容量不断增加，……
\begin{enumerate}
  \item XXXXXXXXXXXXXXXXXXX。
  \item XXXXXXXXXXXXXXXXXXX。
  \item XXXXXXXXXXXXXXXXXXX。
  \item XXXXXXXXXXXXXXXXXXX。
\end{enumerate}

\subsection{拟采取的解决措施}
针对论文工作难点，拟选用……

\section{论文工作量及预期进度}

\begin{schedule}
  \item{20XX年XX月--20XX年XX月}{收集资料，确定设计系统总体方案，翻译有关外文资料及阅读技术文献，书写开题报告。}
  \item{20XX年XX月--20XX年XX月}{学习XXXXXXXXXXXXXX系统开发工具进行设计。}
  \item{20XX年XX月--20XX年XX月}{进行方案论证。}
  \item{20XX年XX月--20XX年XX月}{进行系统总体编程及调试。}
  \item{20XX年XX月--20XX年XX月}{与硬件部分的整合。}
  \item{20XX年XX月--20XX年XX月}{编写毕业论文。}
  \item{20XX年XX月--20XX年XX月}{毕业答辩准备和答辩。}
\end{schedule}

\section{预期成果及其可能的创新点}

预期成果：XXXXXXXXXXXXXXXXXXXXXXXX。

可能的创新点：XXXXXXXXXXXXXXXXXXXXXXXXXX

\nocite{*}
\sdjubibliography

\supervisoropinion{}   % 花括号内可填写指导教师意见
\committeeopinion

\end{document}
